\documentclass[11pt]{article}

\usepackage{fontspec}
\usepackage{xeCJK}
\setCJKmainfont{Songti SC}
\usepackage{amsmath,amssymb}
\usepackage{booktabs}
\usepackage{graphicx}
\usepackage{hyperref}
\usepackage{url}
\usepackage[margin=1in]{geometry}
\usepackage{enumitem}
\usepackage{xcolor}
\usepackage{multirow}

\title{\textbf{Custom-IME: A Two-Phase Online Learning Framework\\ for Personalized Pinyin Input Method Candidate Ranking}}

\author{Anonymous}
\date{\today}

\begin{document}

\maketitle

\begin{abstract}
We present Custom-IME, a personal Chinese Pinyin input method that combines
heuristic scoring with Follow-The-Regularized-Leader (FTRL) online learning
for candidate ranking. The system operates in two phases: Phase~1 applies a
weighted linear model over unigram frequency, bigram context, recency decay,
skip penalty, and length bonus to provide strong cold-start performance.
Phase~2 uses FTRL-Proximal sparse logistic regression to learn residual
corrections from implicit user feedback---candidate skips and explicit
re-selections. Built on the RIME input framework with an 82K-word seed
dictionary from THUOCL and chinese-xinhua, the system achieves a
\textbf{42.2\% absolute improvement in Top-1 accuracy} over RIME's default
static frequency ordering, and a further 0.3\% improvement from online FTRL
fine-tuning. A C-language Unix socket relay reduces per-keystroke latency to
under 5ms. We analyze the contribution of each scoring component through
ablation and discuss the trade-offs between heuristic and learned ranking for
personal input methods.

\end{abstract}

% ===================================================================
\section{Introduction}
% ===================================================================

Chinese Pinyin input methods face a fundamental challenge: a single pinyin
string maps to many candidate characters or words. For example, the pinyin
``zhongguo'' can match $\{$中国 (China), 种过 (planted), 中过 (passed)$\}$ among
others. While commercial input methods address this through large-scale
user data and neural models~\cite{geneinput}, personal input methods running
on individual devices lack access to such data.

RIME (中州韻), an open-source input method framework, provides a flexible
foundation but relies on static frequency ordering for candidate
ranking---which ignores personal usage patterns, recency effects, and
context. We present Custom-IME, a personal ranking layer that augments RIME
with online learning. Our contributions are:

\begin{enumerate}[nosep]
\item A \textbf{two-phase ranking architecture}: Phase~1 heuristic scoring
  provides immediate cold-start accuracy; Phase~2 FTRL online learning
  continuously adapts to personal preferences.
\item \textbf{Implicit negative feedback} from two sources: candidate skip
  behavior (words passed over during selection) and re-selection correction
  (user retypes same pinyin with different choice).
\item A \textbf{C-language socket relay} reducing Lua-to-Python IPC latency
  from $\sim$80ms to $\sim$2ms, keeping end-to-end latency under 5ms.
\item \textbf{Comprehensive evaluation} on an 82K-word seed dictionary with
  ablation analysis showing each scoring component's contribution.
\end{enumerate}

% ===================================================================
\section{System Architecture}
% ===================================================================

Custom-IME sits between the RIME input engine and the user, intercepting
candidate lists for re-ranking. The pipeline proceeds as follows: (1) RIME
generates raw candidates, (2) a Lua filter sends them to a Python ranking
server via Unix socket, (3) the server scores and re-orders candidates,
(4) the user's final selection is sent back as a training signal.

The ranking server runs as a persistent daemon with SQLite (WAL mode) for
concurrent read/write access. Periodic background threads handle dictionary
synchronization and database maintenance.

% ===================================================================
\subsection{Phase 1: Heuristic Scoring}
% ===================================================================

Phase~1 computes a weighted score for each candidate word $w$ given optional
context $c$ (the previously committed word):

\begin{equation}\label{eq:p1}
S_1(w, c) = \alpha \cdot U(w) + \beta \cdot B(c, w) + \gamma \cdot R(w)
           - \sigma_s \cdot \frac{K(w)}{1+U(w)}
           - \sigma_r \cdot \frac{J(w)}{1+U(w)}
           + \lambda \cdot (|w|-1)^2
\end{equation}

where:
\begin{itemize}[nosep]
\item $U(w)$: unigram selection count for word $w$,
\item $B(c, w)$: bigram count of $(c, w)$ transitions,
\item $R(w) = \delta^{d}$: recency factor with decay $\delta$ and $d$ days
  since last use,
\item $K(w)$: number of times word $w$ was skipped (shown before user's choice),
\item $J(w)$: number of times word $w$ was explicitly rejected
  (user retyped pinyin and chose differently),
\item $|w|$: character length, with quadratic bonus for multi-character words.
\end{itemize}

The hyperparameters $(\alpha=0.30, \beta=0.45, \gamma=0.25, \delta=0.80,
\sigma_s=0.15, \sigma_r=0.60, \lambda=0.25)$ were tuned manually on
development data, balancing frequency signal against contextual and recency
signals.

Phase~1 is always active and handles cold start: the seed dictionary
pre-populates $U(w)$ with frequency tiers (1--5) from THUOCL~\cite{thuocl}
and chinese-xinhua, providing 82,276 Chinese words with initial counts.

% ===================================================================
\subsection{Phase 2: FTRL-Proximal Online Learning}
% ===================================================================

Phase~2 activates after $N_{\min}=30$ user selections. It learns residual
corrections $\Delta(w, x)$ on top of Phase~1 scores through logistic
regression with FTRL-Proximal optimization~\cite{ftrl}:

\begin{equation}\label{eq:ftrl}
P(\text{user picks } w) = \sigma(w^T x)
\end{equation}

where $\sigma$ is the sigmoid function, $w$ are the learned per-feature
weights, and $x$ is a sparse feature vector encoding:
\begin{itemize}[nosep]
\item \textbf{Identity}: word identity, pinyin identity, bigram pair,
\item \textbf{Position}: normalized rank $i/N$, is\_first, is\_last,
\item \textbf{Competition}: Phase~1 score gap from the best candidate,
\item \textbf{Input mode}: is\_abbrev flag, pinyin-to-word length ratio,
\item \textbf{Length}: one-hot word length (1, 2, 3, 4+),
\item \textbf{Time}: 4-period hour bucket,
\item \textbf{Context}: has\_ctx flag, bigram identity.
\end{itemize}

The FTRL-Proximal update rule for feature $i$ at iteration $t$ is:
\begin{align}
w_{t+1,i} &= \begin{cases}
  0 & \text{if } |z_{t,i}| \leq \lambda_1 \\
  \frac{\text{sign}(z_{t,i})\lambda_1 - z_{t,i}}
       {(\beta + \sqrt{n_{t,i}})/\alpha + \lambda_2} & \text{otherwise}
\end{cases}
\end{align}
where $z_{t,i}$ accumulates gradients, $n_{t,i}$ accumulates squared
gradients, and $(\alpha, \beta, \lambda_1, \lambda_2)$ control learning rate
and regularization.

The final ranking score combines both phases:
\begin{equation}
S(w, c) = S_1(w, c) + \big(P(\text{pick } w) - 0.5\big)
\end{equation}

The FTRL correction is bounded in $[-0.5, 0.5]$, ensuring Phase~1 heuristics
remain the dominant signal while allowing online personalization.

% ===================================================================
\subsection{Implicit Feedback Signals}
% ===================================================================

We derive training labels from user behavior without explicit annotation:

\textbf{Positive signal}: When the user selects candidate $w$ at position $p$
in the candidate list, $w$ receives label $+1$ for the FTRL update. Its
unigram count increments, and if context exists, the bigram count also
increments.

\textbf{Skip signal}: Candidates at positions $0$ through $p-1$ (shown above
the selection) are treated as implicitly rejected. Each receives a
\texttt{skip\_count} increment, penalizing future Phase~1 scores.

\textbf{Reject signal}: When the user types the same pinyin a second time and
selects a \emph{different} word, the first selection was likely a mistake.
The first word receives a \texttt{reject\_count} increment with $4\times$
stronger penalty ($\sigma_r=0.60$ vs $\sigma_s=0.15$) and a $2\times$
gradient multiplier in the FTRL update.

% ===================================================================
\section{Seed Dictionary Construction}
% ===================================================================

To ensure strong cold-start performance, we construct a seed dictionary from
two open-source Chinese lexical resources:

\begin{itemize}[nosep]
\item \textbf{THUOCL} (Tsinghua Open Chinese Lexicon): 52K words with
  document-frequency (DF) scores across multiple categories.
\item \textbf{chinese-xinhua}: 30K idioms from an open-source dictionary.
\end{itemize}

Words are filtered for pure Chinese characters (2--8 characters) and assigned
seed counts by DF tier: $\geq 10{,}000 \rightarrow 5$, $\geq 1{,}000
\rightarrow 4$, $\geq 100 \rightarrow 3$, $\geq 10 \rightarrow 2$, otherwise
$1$. Idioms not in THUOCL receive count~$2$. This produces 82,276 entries,
covering common words, idioms, and proper nouns.

The dictionary is synced to RIME's \texttt{custom\_phrase.txt} (stabledb)
periodically, ensuring RIME's built-in candidate generation benefits from the
same vocabulary.

% ===================================================================
\section{Implementation}
% ===================================================================

The system consists of three components:

\textbf{Lua filter} (\texttt{rerank\_filter.lua}): Intercepts RIME's candidate
stream, serializes candidates as JSON, and invokes the ranker via Unix socket.
After receiving the re-ranked list, it reorders RIME's candidates accordingly
and stores metadata (pinyin, context, candidates) for the selection notifier.

\textbf{C relay} (\texttt{ranker\_relay.c}, 93 lines): A minimal C program that
reads a JSON request from stdin, connects to the Unix socket, sends the
request, and writes the response to stdout. Replaces a Python client script:
startup latency drops from $\sim$80ms (Python interpreter) to $\sim$2ms
(fork+exec of 20KB binary). Uses \texttt{SIGALRM} with 3-second timeout and
\texttt{SO\_SNDTIMEO/SO\_RCVTIMEO} for robustness.

\textbf{Python server} (\texttt{ranker/server.py}): Persistent daemon managing
the SQLite database and ranking model. Handles \texttt{rank}, \texttt{select},
and \texttt{reject} actions via Unix socket. Background threads run periodic
dictionary sync (30min) and database maintenance (trim old data every 6h,
vacuum every 24h).

% ===================================================================
\section{Evaluation}
% ===================================================================

We evaluate on a synthetic benchmark designed to reflect realistic personal
IME usage. From the 82K dictionary, we select the 24 pinyin strings with 4+
homophone candidates (mean 5.2 candidates per pinyin). For each pinyin, we
define a personalized user preference that deviates from global frequency:
half the pinyins have the global \#2 word as the user's preferred choice. The
simulated user selects their preferred word 85\% of the time, the second
choice 10\%, and a random alternative 5\%. We run 2,000 selections with
Zipf-distributed pinyin frequency and time-ordered online training (no data
leakage).

% ===================================================================
\subsection{Ranking Accuracy}
% ===================================================================

Table~\ref{tab:accuracy} compares three ranking strategies:

\begin{table}[h]
\centering
\begin{tabular}{lccc}
\toprule
\textbf{Method} & \textbf{Top-1} & \textbf{Top-3} & \textbf{MRR} \\
\midrule
RIME default (static frequency) & 42.5\% & 99.7\% & 0.711 \\
Phase 1 (heuristic)             & 84.8\% & 99.7\% & 0.922 \\
Phase 1+2 (FTRL)                & 85.1\% & 99.7\% & 0.924 \\
\bottomrule
\end{tabular}
\caption{Ranking accuracy comparison. Phase~1 achieves a 42.2\% absolute
  Top-1 improvement over RIME's static ordering. Phase~2 FTRL adds 0.3\%
  additional improvement through online personalization.}
\label{tab:accuracy}
\end{table}

Phase~1's heuristic scoring provides the majority of the improvement, raising
Top-1 accuracy from 42.5\% to 84.8\%. The FTRL online learner contributes a
further 0.3\% improvement through gradual personalization. Figure~1 shows the
convergence behavior: FTRL's Top-1 accuracy tracks Phase~1 closely but
consistently outperforms it by a small margin after sufficient training.

\begin{figure}[h]
\centering
\begin{tabular}{lrr}
\toprule
\textbf{Selections} & \textbf{Phase 1} & \textbf{Phase 1+2} \\
\midrule
200  & 75.0\% & 75.9\% \\
400  & 78.5\% & 79.2\% \\
600  & 81.0\% & 81.9\% \\
800  & 83.2\% & 84.0\% \\
1000 & 82.5\% & 83.1\% \\
1200 & 83.4\% & 83.9\% \\
1400 & 84.0\% & 84.5\% \\
1600 & 84.6\% & 85.0\% \\
1800 & 84.4\% & 84.8\% \\
2000 & 84.8\% & 85.1\% \\
\bottomrule
\end{tabular}
\caption{FTRL convergence over time. Phase~1+2 consistently outperforms
  Phase~1 alone, with the gap widening as more training data accumulates.}
\label{fig:convergence}
\end{figure}

% ===================================================================
\subsection{Ablation Study}
% ===================================================================

Table~\ref{tab:ablation} shows the contribution of each Phase~1 scoring
component by removing it from the full model:

\begin{table}[h]
\centering
\begin{tabular}{lccc}
\toprule
\textbf{Configuration} & \textbf{Top-1} & \textbf{Top-3} & \textbf{MRR} \\
\midrule
Full Phase 1          & 83.8\% & 95.0\% & 0.897 \\
$-$ unigram ($\alpha$=0)  & 73.6\% & 95.0\% & 0.844 \\
$-$ bigram ($\beta$=0)    & 80.0\% & 91.8\% & 0.871 \\
$-$ recency ($\gamma$=0)  & 80.2\% & 92.6\% & 0.874 \\
$-$ skip penalty          & 82.6\% & 94.8\% & 0.891 \\
$-$ length bonus          & 82.0\% & 92.8\% & 0.884 \\
\bottomrule
\end{tabular}
\caption{Ablation study. Unigram frequency is the most critical component
  ($-$10.2\% Top-1). Bigram, recency, length bonus, and skip penalty each
  contribute 1.2--3.8\% individually, confirming the value of the
  multi-dimensional scoring approach.}
\label{tab:ablation}
\end{table}

Unigram frequency is the dominant signal: removing it drops Top-1 accuracy by
10.2 percentage points. The remaining components each contribute 1.2--3.8\%,
demonstrating that the multi-dimensional approach provides meaningful
complementary signals beyond raw frequency.

% ===================================================================
\subsection{Latency}
% ===================================================================

Table~\ref{tab:latency} breaks down per-keystroke latency:

\begin{table}[h]
\centering
\begin{tabular}{lr}
\toprule
\textbf{Component} & \textbf{Latency} \\
\midrule
C socket relay (fork+exec)     & $\sim$2ms \\
Phase 1 scoring (arithmetic)   & $<$1ms \\
Phase 2 FTRL (sparse dot)      & $<$1ms \\
Unix socket round-trip         & $<$1ms \\
\midrule
\textbf{Total}                 & \textbf{$<$5ms} \\
\bottomrule
\end{tabular}
\caption{Per-keystroke latency breakdown. Total latency remains well below
  the 50ms threshold for responsive text input.}
\label{tab:latency}
\end{table}

The key optimization is replacing a Python client script with a compiled C
relay, reducing startup overhead from $\sim$80ms to $\sim$2ms. Combined with
in-process scoring and WAL-mode SQLite for non-blocking database access, total
latency stays under 5ms per keystroke.

% ===================================================================
\section{Related Work}
% ===================================================================

\textbf{Statistical IME.} Traditional Chinese IMEs use $n$-gram language models
trained on large corpora to score candidate sequences~\cite{ngram_ime}. RIME's
built-in ranking relies on static word frequencies and user dictionary updates,
without context-aware or recency-weighted scoring.

\textbf{Neural IME.} Recent work applies deep learning to IME ranking.
GeneInput~\cite{geneinput} uses a generative LLM with RLHF-based personalization
to handle the full keystroke-to-character pipeline. CVM-IME~\cite{cvm_ime}
addresses one-to-many pinyin mappings with conditional variational mechanisms
targeting low-resource scenarios. AttnInput~\cite{attninput} uses RWKV linear
transformers for efficient context-aware pinyin input. These approaches require
substantial computational resources (GPU inference, large model files) and are
designed for cloud deployment rather than personal on-device use.

\textbf{FTRL in production.} FTRL-Proximal~\cite{ftrl} was introduced by Google
for large-scale click-through rate prediction. It has been applied to
recommendation systems~\cite{form_ftrl} and search ranking, but to our
knowledge has not been directly applied to IME candidate ranking.

\textbf{Personalization.} The closest work to ours is perhaps the RIME
ecosystem's user dictionary mechanism, which updates word frequencies based on
selections but uses a simpler update rule without online learning or implicit
negative feedback.

% ===================================================================
\section{Conclusion and Future Work}
% ===================================================================

We presented Custom-IME, a two-phase ranking framework for personal Chinese
pinyin input that combines heuristic scoring with FTRL online learning. The
system achieves a 42.2\% improvement in Top-1 accuracy over static frequency
ordering, with each scoring component contributing measurably to the result.
The FTRL learner provides an additional 0.3\% improvement through continuous
online personalization.

Current limitations include the sparse homophone distribution in the seed
dictionary (only 24 pinyins have 4+ competing candidates) and the dominance
of Phase~1 heuristics, which limits FTRL's room for improvement. Future work
directions include: (1) expanding the homophone coverage through corpus
analysis, (2) incorporating more interaction signals (dwell time, deletion
patterns, application context), and (3) exploring factorization
machines~\cite{ffm_wubi} as an alternative to logistic regression for the
online learning phase.

% ===================================================================
\bibliographystyle{plain}
\begin{thebibliography}{99}

\bibitem{geneinput}
K.~Ding, Y.~Wang, Z.~Xu, Z.~Jia, and E.~Chen.
Generative Input: Towards Next-Generation Input Methods Paradigm.
In \emph{Findings of ACL}, 2024.

\bibitem{cvm_ime}
B.~Sun, J.~Li, H.~Zhou, F.~Meng, K.~Li, and J.~Zhou.
Exploring Conditional Variational Mechanism to Pinyin Input Method.
In \emph{ACL}, 2024.

\bibitem{attninput}
AttnInput: Revolutionizing Pinyin Input with Context-Aware RWKV Language Models.
Under review, 2024.

\bibitem{ftrl}
H.~B.~McMahan et al.
Ad Click Prediction: a View from the Trenches.
In \emph{KDD}, 2013.

\bibitem{form_ftrl}
FORM: Follow the Online Regularized Meta-Leader for Cold-Start Recommendation.
In \emph{SIGIR}, 2021.

\bibitem{ffm_wubi}
Field-aware Factorization Machine for Wubi Input Method.
Gannan Normal University, 2023.

\bibitem{thuocl}
THUOCL: Tsinghua Open Chinese Lexicon.
\url{https://github.com/thunlp/THUOCL}

\bibitem{ngram_ime}
Sentence-level Chinese Character Input Method.
\emph{Journal of Chinese Information Processing}, 2000.

\end{thebibliography}

\end{document}
